\documentclass{article}
\usepackage[margin=1in, a4paper]{geometry}
\usepackage{amsmath, amssymb, amsfonts}
\usepackage{graphicx}
\usepackage{xcolor}
\usepackage{listings}
\usepackage{booktabs}
\usepackage[utf8]{inputenc}
\usepackage[T1]{fontenc}
\usepackage{hyperref}
\hypersetup{colorlinks=true, urlcolor=blue, linkcolor=blue}
\usepackage{enumitem}
\usepackage{float}

\definecolor{codegreen}{rgb}{0,0.6,0}
\definecolor{codegray}{rgb}{0.5,0.5,0.5}
\definecolor{codepurple}{rgb}{0.58,0,0.82}
\definecolor{backcolour}{rgb}{0.99,0.99,0.99}
% Professional color palette for academic publishing
\definecolor{myblue}{cmyk}{0.8,0.4,0,0.1}
\definecolor{mygreen}{cmyk}{0.7,0,0.5,0.2}
\definecolor{myorange}{cmyk}{0,0.5,0.8,0.1}
\definecolor{mygray}{cmyk}{0,0,0,0.4}
\definecolor{arrowcolor}{cmyk}{0.9,0.5,0,0.3}
\definecolor{nodecolor}{cmyk}{0.6,0.2,0,0.05}
\definecolor{framecolor}{cmyk}{0.4,0.1,0,0.1}

\lstdefinestyle{mystyle}{
    backgroundcolor=\color{backcolour},   
    commentstyle=\color{codegreen},
    keywordstyle=\color{magenta},
    numberstyle=\tiny\color{codegray},
    stringstyle=\color{codepurple},
    basicstyle=\ttfamily\footnotesize,
    breakatwhitespace=false,         
    breaklines=true,                 
    captionpos=b,                    
    keepspaces=true,                 
    numbers=left,                    
    numbersep=5pt,                  
    showspaces=false,                
    showstringspaces=false,
    showtabs=false,                  
    tabsize=2
}
\lstset{style=mystyle}

\title{The New Product Forecasting Problem (Look-Alike Modeling)}
\author{The Logistics \& Supply Chain Problem-Solving Playbook}
\date{}

\begin{document}

\maketitle

\section{The New Product Forecasting Problem}

New product forecasting represents one of the most challenging aspects of supply chain management, particularly when historical sales data is limited or non-existent. Traditional forecasting methods rely heavily on past performance patterns, making them ineffective for newly launched products. Look-alike modeling emerges as a sophisticated solution that leverages similarities between new products and existing ones to generate reliable demand forecasts, enabling organizations to optimize inventory levels, production planning, and market positioning strategies.

\subsection{The Scenario: Launching the Next Generation Smartphone}

TechNova Corporation, a leading electronics manufacturer, faces a critical challenge as they prepare to launch their revolutionary smartphone model, the \texttt{TechNova\_X1}. With an anticipated market introduction in Q2 2026, the company must forecast demand across multiple geographic markets and retail channels to optimize their supply chain operations. The product features cutting-edge technology including advanced AI capabilities, extended battery life, and innovative design elements that differentiate it significantly from their existing product portfolio.

The challenge lies in accurately predicting demand patterns when no direct historical data exists for this specific product. Traditional time-series forecasting methods prove inadequate, as they require substantial historical sales data. The company's supply chain team must determine optimal production quantities, distribution strategies, and inventory allocation across their global network of manufacturing facilities and distribution centers.

The stakes are substantial: overproduction could result in millions of dollars in excess inventory and storage costs, while underproduction could lead to stockouts, lost sales, and damaged market reputation. TechNova's management recognizes that successful demand forecasting will directly impact their competitive positioning and financial performance in the highly competitive smartphone market[2].

\subsection{Tier 1: The Pen \& Paper Method (Mathematical Formulation)}

The mathematical foundation for new product forecasting using look-alike modeling centers on a scheduling optimization framework that allocates forecasting resources across time periods while minimizing prediction errors and meeting market coverage requirements.

\textbf{Mathematical Model Formulation}

Let us define the following sets and parameters:
\begin{align}
P &= \{1, 2, \ldots, n\} \quad \text{Set of existing products (potential analogs)} \\
T &= \{1, 2, \ldots, m\} \quad \text{Set of time periods for forecasting} \\
S &= \{1, 2, \ldots, k\} \quad \text{Set of market segments} \\
s_{it} &= \text{Sales of product } i \text{ in period } t \\
\alpha_{ij} &= \text{Similarity coefficient between products } i \text{ and } j \\
d_{jt} &= \text{Forecasted demand for new product } j \text{ in period } t \\
w_i &= \text{Weight assigned to product } i \text{ as analog} \\
C &= \text{Maximum number of analog products to consider}
\end{align}

The objective function minimizes the total forecasting error while balancing computational complexity:

\begin{equation}
\text{Minimize} \quad Z = \sum_{t \in T} \sum_{s \in S} \left( d_{jt}^s - \sum_{i \in P} w_i \cdot s_{it}^s \cdot \alpha_{ij} \right)^2 + \lambda \sum_{i \in P} w_i
\end{equation}

Subject to the following constraints:

\textbf{Analog Selection Constraint:}
\begin{equation}
\sum_{i \in P} x_i \leq C
\end{equation}

\textbf{Weight Normalization Constraint:}
\begin{equation}
\sum_{i \in P} w_i = 1
\end{equation}

\textbf{Similarity Threshold Constraint:}
\begin{equation}
w_i \leq M \cdot x_i \quad \forall i \in P
\end{equation}

\textbf{Minimum Similarity Constraint:}
\begin{equation}
\alpha_{ij} \cdot x_i \geq \theta \cdot x_i \quad \forall i \in P
\end{equation}

\textbf{Demand Non-negativity Constraint:}
\begin{equation}
d_{jt}^s \geq 0 \quad \forall t \in T, s \in S
\end{equation}

Where $x_i \in \{0,1\}$ indicates whether product $i$ is selected as an analog, $M$ is a large positive constant, $\lambda$ is the regularization parameter, and $\theta$ is the minimum similarity threshold.

\begin{figure}[htbp]
\centering
\includegraphics[width=\textwidth]{../figures-png/line50.fig1.png}
\caption{Enhanced New Product Development Problem with Professional CMYK Colors and Node-Based Similarity Network}
\end{figure}

\textbf{Concrete Mathematical Example}

Consider TechNova's scenario with 4 existing smartphone models and 3 forecast periods. The similarity coefficients matrix and historical sales data are:

Similarity coefficients: $\alpha = [0.85, 0.72, 0.64, 0.42]$

Historical sales data (units in thousands):
\begin{align}
\mathbf{S} = \begin{pmatrix}
120 & 95 & 78 \\
85 & 110 & 102 \\
95 & 88 & 93 \\
65 & 72 & 68
\end{pmatrix}
\end{align}

Solving the optimization problem with $C = 2$ (maximum 2 analogs), we obtain optimal weights: $w = [0.65, 0.35, 0, 0]$, selecting products 1 and 2 as the primary analogs.

The forecasted demand for the new product becomes:
\begin{align}
d_{1} &= 0.65 \times 120 \times 0.85 + 0.35 \times 85 \times 0.72 = 87.6 \text{ thousand units} \\
d_{2} &= 0.65 \times 95 \times 0.85 + 0.35 \times 110 \times 0.72 = 80.1 \text{ thousand units} \\
d_{3} &= 0.65 \times 78 \times 0.85 + 0.35 \times 102 \times 0.72 = 68.4 \text{ thousand units}
\end{align}

\subsection{Tier 2: The Classic Heuristic (Python Implementation)}

The heuristic approach employs Lagrangian relaxation to decompose the complex look-alike modeling problem into manageable subproblems, enabling efficient solution of large-scale forecasting instances without requiring sophisticated optimization solvers.

\textbf{Algorithm Conceptual Framework}

The Lagrangian relaxation heuristic works by relaxing the coupling constraints that link product selection decisions across different market segments and time periods. This decomposition allows us to solve smaller, independent subproblems for each segment-period combination, then coordinate the solutions through an iterative price adjustment mechanism.

\begin{figure}[htbp]
\centering
\includegraphics[width=\textwidth]{../figures-png/line50.fig2.png}
\caption{Lagrangian Relaxation Decomposition for Look-Alike Modeling}
\end{figure}

\textbf{Pseudocode Implementation}

\lstinputlisting[language=Python, caption=Lagrangian Relaxation Heuristic for Look-Alike Modeling]{.././codes-python/line50.py1.py}

\textbf{Concrete Heuristic Example}

Applying the Lagrangian relaxation heuristic to TechNova's smartphone forecasting problem with 3 market segments (Premium, Mid-range, Budget) and 4 existing products:

\textbf{Input Data:}
- Products: [iPhone\_Pro, Galaxy\_Ultra, Pixel\_Advanced, OnePlus\_Elite]
- Similarity scores to new product: [0.85, 0.72, 0.64, 0.58]
- Market segments: 3 (Premium, Mid-range, Budget)
- Time periods: 3 (Launch, Growth, Maturity)
- Maximum analogs: 2

\textbf{Algorithm Execution Results:}

Iteration 1: Initial multipliers $\lambda_w = [0.1, 0.1, 0.1]$, $\lambda_s = 0.05$
- Segment solutions: Premium selects [iPhone\_Pro, Galaxy\_Ultra], Mid-range selects [Galaxy\_Ultra, Pixel\_Advanced], Budget selects [Pixel\_Advanced, OnePlus\_Elite]
- Constraint violation: 0.23 (weight normalization across segments)

Iteration 15: Converged multipliers $\lambda_w = [0.087, 0.094, 0.089]$, $\lambda_s = 0.041$
- Coordinated solution: Primary analogs [iPhone\_Pro, Galaxy\_Ultra] with weights [0.58, 0.42]
- Final forecast: Premium segment: 89.2K, 82.1K, 74.3K units across 3 periods

The heuristic achieves 94.2\% accuracy compared to optimal solution while reducing computation time from 47 minutes to 3.2 minutes for the full problem instance[1].

\subsection{Tier 3: The Advanced Algorithm (Metaheuristic Implementation)}

The Krill Herd Algorithm represents a sophisticated nature-inspired metaheuristic that models the complex herding behavior of Antarctic krill to solve the multi-dimensional look-alike modeling optimization problem. This approach excels in navigating the complex solution space where traditional methods struggle with local optima.

\textbf{Krill Herd Algorithm Theory}

The Krill Herd Algorithm mimics the foraging and herding behavior of krill in marine ecosystems. Each krill individual represents a potential solution combining product selection, weight allocation, and forecasting parameters. The algorithm balances three key behavioral components: movement induced by other krill (social behavior), foraging activity (exploitation), and random diffusion (exploration).

The position of the $i$-th krill is updated according to:
\begin{equation}
\frac{dx_i}{dt} = N_i + F_i + D_i
\end{equation}

Where $N_i$ represents the motion induced by neighbors, $F_i$ represents the foraging motion, and $D_i$ represents the random diffusion motion.

\begin{figure}[htbp]
\centering
\includegraphics[width=\textwidth]{../figures-png/line50.fig3.png}
\caption{Krill Herd Algorithm Behavior in Look-Alike Solution Space}
\end{figure}

\textbf{Pseudocode Implementation}

\lstinputlisting[language=Python, caption=Krill Herd Algorithm for Look-Alike Modeling]{.././codes-python/line50.py2.py}

\textbf{Concrete Metaheuristic Example}

Testing the Krill Herd Algorithm on TechNova's expanded dataset with 8 historical smartphone models across 4 market segments and 6 forecast periods:

\textbf{Algorithm Configuration:}
- Population size: 50 krill individuals
- Maximum iterations: 200
- Inertia weights: [Motion: 0.9, Foraging: 0.2, Diffusion: 0.1]
- Crossover rate: 0.8, Mutation rate: 0.05

\textbf{Performance Results:}

Initial Population (Iteration 0):
- Best fitness: 1,247.3 (high forecasting error)
- Selected analogs: Random selection of 4-6 products
- Average forecast accuracy: 62.4\%

Mid-Evolution (Iteration 100):
- Best fitness: 156.8 (significant improvement)
- Convergent analog selection: [iPhone\_Pro, Galaxy\_Ultra, Pixel\_Advanced]
- Improved forecast accuracy: 87.6\%

Final Solution (Iteration 187):
- Best fitness: 89.2 (near-optimal solution)
- Optimized analog selection: [iPhone\_Pro, Galaxy\_Ultra] with weights [0.62, 0.38]
- Final forecast accuracy: 93.8\%
- Computational time: 12.4 minutes

The Krill Herd Algorithm demonstrates superior exploration capabilities, discovering high-quality solutions that traditional optimization methods miss. The marine ecosystem metaphor effectively balances exploitation of promising regions with exploration of diverse solution alternatives, achieving 93.8\% forecast accuracy compared to 89.1\% from conventional approaches[1].

\subsection{Tier 4: The AI/ML/RL Augmentation Method}

The adversarial training framework enhances look-alike modeling robustness by training forecasting models against adversarial perturbations in product similarity data, market conditions, and historical patterns. This approach ensures reliable performance under uncertain and dynamically changing market environments.

\textbf{Adversarial Training Framework Theory}

Adversarial training works by simultaneously optimizing two competing objectives: a forecasting model that aims to minimize prediction errors, and an adversarial model that generates perturbations designed to maximize those same errors. This minimax optimization creates robust models that maintain accuracy even when input data contains noise, outliers, or systematic biases.

The mathematical formulation follows:
\begin{align}
\min_{\theta} \max_{\delta} \mathbb{E}_{(x,y) \sim D} [L(f_\theta(x + \delta), y)]
\end{align}

Where $f_\theta$ represents the forecasting model with parameters $\theta$, $\delta$ represents adversarial perturbations, $L$ is the loss function, and $D$ is the training data distribution.

\begin{figure}[htbp]
\centering
\includegraphics[width=\textwidth]{../figures-png/line50.fig4.png}
\caption{Adversarial Training Architecture for Robust Look-Alike Forecasting}
\end{figure}

\textbf{Implementation Framework}

\lstinputlisting[language=Python, caption=Adversarial Training for Look-Alike Forecasting]{.././codes-python/line50.py3.py}

\textbf{Concrete AI/ML Example}

Implementing adversarial training for TechNova's smartphone forecasting with comprehensive feature engineering:

\textbf{Dataset Configuration:}
- Historical products: 12 smartphone models (2019-2025)
- Feature dimensions: 156 (product attributes, market context, temporal features)
- Market segments: 5 (Premium, Upper-mid, Mid-range, Budget, Entry-level)
- Forecast horizon: 6 periods (months)

\textbf{Training Configuration:}
- Epochs: 200
- Batch size: 32
- Adversarial perturbation budget: $\epsilon = 0.1$
- Attack methods: FGSM, PGD-7, Learned adversarial

\textbf{Performance Results:}

Standard Neural Network Baseline:
- Clean data accuracy: 91.2\%
- Adversarial robustness (FGSM): 67.3\%
- Adversarial robustness (PGD): 52.8\%
- Mean absolute error: 8,247 units

Adversarial Trained Model:
- Clean data accuracy: 89.6\% (slight decrease)
- Adversarial robustness (FGSM): 86.4\% (significant improvement)
- Adversarial robustness (PGD): 82.1\% (significant improvement)
- Mean absolute error: 8,956 units
- Robustness score: 0.847

\textbf{Uncertainty Quantification Results:}

For TechNova X1 launch forecast:
- Point prediction: 142,350 units (Month 1)
- 95\% confidence interval: [138,220 - 146,480] units
- Prediction uncertainty: ±2.9\%
- Robustness under market volatility: 84.7\% maintained accuracy

The adversarial training framework provides robust forecasting capabilities that maintain high accuracy even when market conditions deviate significantly from training scenarios, crucial for new product launches in uncertain environments[2].

\subsection{Tier 7: The Human-AI Symbiotic Partnership (The Centaur Model)}

The Human-AI Symbiotic Partnership transforms look-alike modeling from a purely algorithmic process into a collaborative intelligence system where human domain expertise and AI computational power combine to create superior forecasting capabilities. This approach recognizes that new product forecasting requires both pattern recognition and contextual understanding that neither humans nor AI can achieve optimally in isolation.

\textbf{Symbiotic Architecture Framework}

The centaur model integrates human expertise at critical decision points while leveraging AI for data processing, pattern detection, and computational optimization. The system creates a continuous feedback loop where human insights inform AI model parameters, and AI outputs provide data-driven evidence for human strategic decisions.

The partnership operates through multiple interaction modalities:

\textbf{Expert-Guided Feature Engineering:} Domain experts identify relevant product attributes, market dynamics, and contextual factors that algorithms might miss. For instance, experts might recognize that ``smartphone camera quality'' strongly correlates with ``social media engagement levels'' in certain demographics, leading to enhanced similarity metrics.

\textbf{Interactive Model Calibration:} Rather than treating model parameters as fixed algorithmic outputs, the system allows experts to interactively adjust similarity thresholds, analog selection criteria, and forecast adjustments based on market intelligence that quantitative data cannot capture.

\textbf{Contextual Override Mechanisms:} The system provides structured interfaces for experts to override algorithmic recommendations when contextual factors (regulatory changes, competitive dynamics, seasonal variations) suggest alternative forecasting approaches.

\begin{figure}[htbp]
\centering
\includegraphics[width=\textwidth]{../figures-png/line50.fig5.png}
\caption{Human-AI Symbiotic Partnership Architecture for Look-Alike Forecasting}
\end{figure}

\textbf{Implementation Framework}

\lstinputlisting[language=Python, caption=Human-AI Symbiotic Look-Alike Forecasting System]{.././codes-python/line50.py4.py}

\textbf{Concrete Human-AI Partnership Example}

Implementing the symbiotic system for TechNova's X1 smartphone launch:

\textbf{Expert Team Configuration:}
- Lead Product Manager (15 years smartphone industry experience)
- Market Research Director (expertise in consumer electronics trends)  
- Supply Chain Analytics Expert (forecasting domain specialist)
- Regional Sales Managers (4 geographic markets)

\textbf{Collaborative Forecasting Session Results:}

\textbf{Phase 1: Enhanced Similarity Assessment}
- AI Algorithm: iPhone\_Pro (0.87), Galaxy\_Ultra (0.79), Pixel\_Advanced (0.71)
- Expert Adjustments: 
  - iPhone\_Pro adjusted to 0.82 (``Different pricing strategy reduces similarity'')
  - Galaxy\_Ultra adjusted to 0.84 (``Similar premium positioning and features'')
  - OnePlus\_Elite added at 0.76 (``Overlooked by algorithm - similar target demographic'')

\textbf{Phase 2: Analog Selection Validation}
- AI Recommendation: [iPhone\_Pro, Galaxy\_Ultra, Pixel\_Advanced]
- Expert Validation: [Galaxy\_Ultra, OnePlus\_Elite, iPhone\_Pro]
- Reasoning: ``Galaxy\_Ultra prioritized due to similar Android ecosystem and premium positioning''

\textbf{Phase 3: Contextual Forecast Generation}

Q2 2026 Launch Period:
- Base AI Forecast: 148,500 units
- Expert Market Context Adjustments:
  - Seasonal factor: 1.05 (``Spring smartphone upgrade cycle'')
  - Competitive pressure: 1.15 (``Apple iPhone refresh expected in Q3'')
  - Economic conditions: 0.97 (``Consumer spending uncertainty'')
  - Technology adoption: 1.22 (``High enthusiasm for AI features'')
- Final Symbiotic Forecast: 164,230 units

\textbf{Performance Comparison (6-month validation):}

Traditional AI-Only Forecasting:
- Average forecast error: 12.4\%
- Prediction confidence: 78\%
- Time to generate forecast: 15 minutes

Symbiotic Human-AI System:
- Average forecast error: 6.8\% (45\% improvement)
- Prediction confidence: 91\% (expert validation increases confidence)
- Time to generate forecast: 47 minutes (includes expert consultation)
- Explainability score: 94\% (high stakeholder understanding)

\textbf{Key Success Factors:}

The symbiotic approach demonstrates superior performance through several mechanisms. Expert domain knowledge compensates for algorithmic blind spots, particularly in identifying market context factors that historical data cannot capture. The collaborative process creates more robust forecasts by combining pattern recognition capabilities of AI with contextual understanding of human experts. Most importantly, the system generates highly explainable forecasts that stakeholders trust and can act upon with confidence, leading to better supply chain decisions and reduced inventory risks[3][4].

The continuous learning mechanism adapts collaboration weights based on forecast performance, creating an evolving partnership that becomes more effective over time. This approach represents a paradigm shift from replacing human expertise with AI to augmenting human capabilities through intelligent collaboration.

## Practice Questions

\textbf{Tier 1 Questions (Mathematical Formulation):}

\textbf{Question 1:} A consumer electronics company is launching a new tablet model and has identified 5 potential analog products with the following similarity coefficients and 3-period historical sales data (in thousands):

Product similarities: $\alpha = [0.82, 0.76, 0.69, 0.58, 0.45]$

Historical sales matrix:
$\mathbf{S} = \begin{pmatrix} 95 & 87 & 79 \\ 110 & 102 & 94 \\ 78 & 85 & 88 \\ 125 & 118 & 112 \\ 65 & 71 & 68 \end{pmatrix}$

Formulate the complete mathematical optimization model for selecting at most 2 analog products and determining optimal weights. Calculate the forecasted demand for all 3 periods if the optimal solution selects products 2 and 4 with weights $w_2 = 0.6$ and $w_4 = 0.4$.

\textbf{Question 2:} Extend the mathematical model to include market segment differentiation. Given 2 market segments (Premium and Standard) with different similarity thresholds ($\theta_{premium} = 0.75$, $\theta_{standard} = 0.60$) and segment-specific sales data, formulate the multi-segment optimization problem and determine the optimal analog selection strategy.

\textbf{Tier 2 Questions (Heuristic Algorithms):}

\textbf{Question 3:} Implement the Lagrangian relaxation heuristic for a look-alike modeling problem with 8 existing products, 4 market segments, and 5 time periods. Given initial Lagrangian multipliers $\lambda_w = [0.12, 0.08, 0.15, 0.10]$ and $\lambda_s = 0.06$, describe the decomposition strategy and calculate the first iteration subproblem solutions. How would you coordinate the segment-specific solutions to ensure global feasibility?

\textbf{Question 4:} Design a priority-based heuristic algorithm for real-time analog selection in an e-commerce environment where new products are continuously introduced. The algorithm must handle 1000+ existing products and provide analog recommendations within 2 seconds. Specify the data structures, priority scoring mechanism, and computational complexity analysis.

\textbf{Tier 3 Questions (Metaheuristic Algorithms):}

\textbf{Question 5:} Configure a Krill Herd Algorithm for look-alike modeling with the following parameters: 40 krill individuals, similarity search space [0,1]$^{12}$, weight allocation space simplex constraints. Design the fitness function incorporating forecast accuracy, model complexity, and solution diversity. Describe the neighbor interaction, foraging behavior, and random diffusion mechanisms specific to this optimization landscape.

\textbf{Question 6:} Compare the performance of Krill Herd Algorithm against Particle Swarm Optimization and Genetic Algorithm for a large-scale look-alike modeling problem (50 products, 8 segments, 12 periods). Design experiments to evaluate convergence speed, solution quality, and robustness to local optima. What hybrid approaches could combine the strengths of multiple metaheuristics?

\textbf{Tier 4 Questions (AI/ML/RL Methods):}

\textbf{Question 7:} Implement adversarial training for robust look-alike forecasting where the model must maintain accuracy under input perturbations up to $\epsilon = 0.15$. Given a neural network with architecture [128, 64, 32, 6] and training dataset of 50,000 product-forecast pairs, design the FGSM and PGD attack strategies. Calculate the expected robustness improvement and computational overhead compared to standard training.

\textbf{Question 8:} Design a multi-task learning framework that simultaneously performs look-alike modeling, demand forecasting, and price optimization for new product launches. Specify the shared representation layers, task-specific output heads, and loss function weighting strategy. How would you handle the different scales and objectives across the three tasks?

\textbf{Tier 7 Questions (Human-AI Symbiotic Partnership):}

\textbf{Question 9:} Design a collaborative interface for domain experts to interact with the look-alike forecasting AI system. The interface must capture expert similarity assessments, analog validation decisions, and market context insights while maintaining efficiency for busy executives. Specify the interaction modalities, explanation generation mechanisms, and trust calibration strategies.

\textbf{Question 10:} Implement an adaptive collaboration learning system that adjusts human-AI partnership weights based on forecast performance over time. Given historical performance data showing 15\% average error for AI-only forecasts and 8\% error for expert-guided forecasts, design the weight update mechanism. How would you handle scenarios where expert availability varies or multiple experts provide conflicting insights?

\end{document}
